% SPDX-License-Identifier: Apache-2.0
\appendix
\section{Statusabgleich und Grenzen der Aussagen}
\label{sec:implementation}
Diese Fassung wurde am 29. Oktober 2026 gegen den lokalen Arbeitsstand
auf Basis des Commits \code{8589c5e499cc} geprüft. Die begleitenden
Dokumentationskorrekturen gehören zu dieser Lieferung; der Commitbezug
allein enthält sie noch nicht. Die folgenden Befunde beziehen sich auf
gelesene lokale Implementierungen, nicht auf eine verifizierte
Produktivinstallation oder eine vollständige Konformitätsprüfung.

\begin{longtable}{@{}P{32mm}P{116mm}@{}}
\toprule
Gegenstand & Geprüfter Unterschied und Konsequenz\\
\midrule\endhead
Clustererfüllung & Die normative Formel berücksichtigt alle atomaren
Nachfahren. \code{LearnerService.computeEffectivePrereqMastery} und
\code{computeScopedPrerequisiteMastery} bevorzugen dagegen rekursiv
Kinder, die \code{isCoreForPrereqs} als Kern einstuft, sobald ein
solches Kind existiert. Diese Einstufung verwendet dort fehlende Tags
beziehungsweise \code{GK}-Tags, nicht einfach das Feld \code{core}.
Beide Regeln sind nicht identisch.\\
Binärer Zielzustand & Verbindlich ist geschafft oder noch nicht geschafft.
Die aktuelle \code{Mastery}-Entität und der Schreibpfad akzeptieren noch
numerische Zwischenwerte; \code{GoalCard} bietet einen Regler und
0/50/100-Schaltflächen. Dienste verwenden teilweise die Schwelle
$0{,}9$. Diese Stellen entsprechen nicht dem binären Sollmodell.
Eine numerische Kodierung durch ausschließlich 0 und 1 wäre möglich;
Zwischenwerte sind keine weiteren zulässigen Ergebniszustände.\\
Gedächtnisabschluss & Der dokumentierte Zieltypvertrag fordert
verifizierten Abruf und keine fälligen Karten. Backend und
\code{useSrsMastery} verwenden derzeit allein den verifizierten Abruf
aller erforderlichen Karten. Erneute Fälligkeit und erneute
Fehlprüfung haben daher unterschiedliche Auswirkungen.\\
Graphvalidierung & Die Spezifikation fordert lokale und transitive
Minimalität. Der dokumentierte und gelesene Core-Validator weist keine
vollständige entsprechende Minimalitätsprüfung aus. Ein grüner
Strukturcheck ist deshalb kein vollständiger Konformitätsbeweis.\\
Curriculare Level & Die mathematischen Machbarkeitsbedingungen sind ein
optionales Konzeptprofil. Es wird kein produktiv aktivierter neuer
Levelresolver oder Budgetvalidator behauptet.\\
Prüfungsabschluss & Direkter Prüfzugang und Freistellungsfixpunkt sind
ein Entwurf. Er verwendet vorhandene Mastery als einzige individuelle
Ergebniseingabe; der heutige Ablauf berechnet daraus noch keinen
solchen Umfangsabschluss.\\
\bottomrule
\end{longtable}

Die Ausarbeitung ändert keinen dieser Laufzeitpfade. Eine Angleichung
benötigt eine ausdrückliche Entscheidung zum beabsichtigten Verhalten,
geeignete Regressionen und gegebenenfalls einen Umgang mit vorhandenen
Lernergebnissen. Die Unterschiede als offene Modell-Implementierungsfragen
zu benennen ist genauer, als sie durch eine stillschweigende neue Formel
zu verdecken.

\begin{center}\begin{minipage}{\linewidth}
\paragraph{Prüfbare Codeanker.}
\srcref{backend/src/main/java/com/skillpilot/backend/service/LearnerService.java}{LearnerService.java}:
\begin{itemize}\small
\item \code{isSrsMasteredByVerifiedRecall}
\item \code{computeEffectivePrereqMastery}
\item \code{computeScopedPrerequisiteMastery}
\item \code{isCoreForPrereqs}
\end{itemize}
\srcref{app/src/hooks/useSrsMastery.ts}{useSrsMastery.ts};
\srcref{app/scripts/validateGraph.ts}{validateGraph.ts};
\srcref{backend/src/main/java/com/skillpilot/backend/domain/Mastery.java}{Mastery.java};
\srcref{app/src/components/GoalCard.tsx}{GoalCard.tsx};
\srcref{docs/qa-ci/graph-validation-rules.md}{Graph Validation Rules}.
Die Links zeigen auf die weiterentwickelte Hauptlinie; für eine historische
Prüfung ist zusätzlich der genannte Stand zu verwenden.
\end{minipage}\end{center}

\section{Prüfstrategie für den Abschlussentwurf}
Das begleitende \code{latex/check-assessment-model.py} ist eine
ausführbare mathematische Referenzprüfung mit ausschließlich
Python-Standardbibliothek. Es liest keine Lernerdaten und verändert
keinen Runtimezustand. Es ergänzt die allgemeinen Beweise durch
systematische endliche Gegenbeispielsuche.

\begin{center}\begin{tabularx}{\linewidth}{@{}P{40mm}X@{}}
\toprule
Prüfbereich & Geprüfte Aussage\\
\midrule
Konkrete Fälle & Geteilte Grundlagen, indirekte Verbraucher,
unterstützende Nichtziele, spätere externe Verbraucher,
geschützte Ziele, ungültige Referenzen und erreichbare Zyklen.\\
Mastery als einzige Eingabe & Vorhandene Prüfungs-Mastery zählt;
$E=M_A\cap V_H$; externe Mastery verändert den Kontext nicht;
Freistellungen verändern keine Mastery und werden bei einer
Mastery-Korrektur neu berechnet.\\
Alle beschrifteten DAGs bis vier Knoten & Fixpunktgleichheit mit
topologischer Berechnung; kleinster Fixpunkt; Rundengrenze;
Schutz; keine offenen Verbraucher eines freigestellten Ziels;
Monotonie bei wachsender Mastery-Menge.\\
Regelverhalten bis drei Knoten & Nur Mastery der erforderlichen
Prüfungsziele aktiviert geprüfte Regeln; zusätzliche Mastery erhält
Abschluss oder Freistellung.\\
Kontextgrenzen & Eine unabhängige atomare Expansionsschleife
blockiert keine fremde lokale Berechnung; ein vom Ziel erreichbarer
Zyklus wird abgelehnt.\\
\bottomrule
\end{tabularx}\end{center}

Diese Prüfung enthält keine Modelle realer Aufgabenqualität,
Benotungszuverlässigkeit, Identitätskontrolle oder pädagogischer Wirkung.
Eine spätere Umsetzung muss zusätzlich konkurrierende Änderungen,
idempotente Mastery-Speicherung, Rechte, korrigierte Mastery und
übereinstimmende Anzeigen prüfen. Der Referenzalgorithmus
ist kein neuer produktiver Service.

\section{Begriffe und Symbole}
\begin{longtable}{@{}P{36mm}P{112mm}@{}}
\toprule
Zeichen / Begriff & Bedeutung\\
\midrule\endhead
$G,A,K$ & Ziele, strukturelle Atome, Cluster.\\
$C,R,\Reff$ & Enthaltensein, direkte Voraussetzungen,
direkte plus geerbte Voraussetzungen.\\
$X^+,X^{-1}$ & Nichtleere transitive Hülle und Umkehrrelation.\\
$\Anc,\Desc$ & Vorfahren und Nachfahren; in der Quellspezifikation
\code{Ancestors} und \code{Descendants}.\\
$\Atoms(g)$ & Atomare Basis; bei einem Atom die einelementige Menge.\\
$O,Z,E_{\mathrm{route}}$ & Orientierungsanker, autonome Endziele und
explizite Routenausnahmen; $O,Z$ entsprechen dort $M,T$.\\
$R_A$ & Direkte atomare Voraussetzungskanten $R\cap(A\times A)$.\\
$M_A,m,\Sat$ & Geschaffte Atome, ihr binärer Indikator
$m(a)\in\{0,1\}$ und daraus abgeleitete Erfüllung eines Ziels.\\
$F,G_F$ & Filterprädikat und ausgewählter Zielbereich.\\
$D,V_d$ & Filterdimensionen und Werte einer Dimension $d$.\\
Frontier & Noch nicht erfüllte atomare Ziele mit erfüllten
Voraussetzungen im gewählten Prüfmodus.\\
Personal Curriculum & Verbindlicher längerfristiger persönlicher
Zielumfang; Personalisierungslevel 2.\\
Fokus & Vorübergehende Arbeitsauswahl im Personal Curriculum;
Personalisierungslevel 3.\\
$T_\pi,B_\pi,\ell_\pi$ & Zielmenge, Eintrittsbasis und Erwerbslevel
eines curricularen Plans.\\
$T_i,Q_i$ & Ziele und erforderliche Prüfungsziele eines Abschlussprofils.\\
$F_{\mathrm{atom}},D_{\mathrm{atom}}$ & Atomar expandierte Anforderungen
und deren transitive Hülle; im Abschluss-Quelldokument $F,D$.\\
$H=(V_H,R_H)$ & Unter atomaren Voraussetzungen abgeschlossener
Berechnungskontext mit transitiver Bedarfsrelation.\\
$E,P$ & $E=M_A\cap V_H$: vorhandene Mastery im Kontext;
$P$ enthält geschützte Ziele.\\
$S_r,W_r,W(E)$ & Zu schaffende Prüfungsziele einer Regel, deren mögliche
Freistellungen und aktuell aktivierte Freistellungsberechtigung.\\
$X_*,I$ & Berechnete erledigte Arbeit durch Mastery oder Freistellung;
$I=X_*\setminus E$ enthält nur Freistellungen. Keine zusätzlichen
gespeicherten Ergebniszustände.\\
M7 & Benannter maschineller Curriculum-QA-Reifegrad,
kein menschlicher Wirksamkeitsnachweis.\\
\bottomrule
\end{longtable}

\section{Quellenzuordnung und Abbildungen}
Die Kapitel sind eine redaktionell zusammengeführte deutsche Darstellung
der folgenden Repositoryquellen. Querverweise nennen die jeweilige
fachliche Verantwortung; sie sind keine Behauptung, dass alle optionalen
Konzepte bereits implementiert seien.

\begin{longtable}{@{}P{38mm}P{110mm}@{}}
\toprule
Inhalt & Maßgebliche Quelle\\
\midrule\endhead
Formaler Kern und Filter &
\srcref{docs/concept/skill-graph/graph-definition.md}{Graph Definition},
Abschnitte 1--11.\\
Schichten und Migration &
\srcref{docs/concept/skill-graph/general-goal-system-and-migration.md}{General Goal System and Migration};
\srcref{docs/concept/skill-graph/canonical-gymnasium-rollout.md}{Canonical Gymnasium Rollout}.\\
Sichten und persönliche Auswahl &
\srcref{docs/concept/skill-graph/view-projection-and-goal-placement.md}{View Projection and Goal Placement};
\srcref{docs/concept/levels-of-personalization.md}{Levels of Personalization}.\\
Zieltypen und Evidenz &
\srcref{docs/concept/skill-graph/node-types.md}{Node Types};
\srcref{AGENTS.md}{AGENTS.md, insbesondere Orientierung und Mastery}.\\
Curriculare Planung &
\srcref{docs/concept/skill-graph/graph-definition.md}{Graph Definition, Abschnitt 12};
\srcref{docs/concept/skill-graph/curriculum-levels.md}{Curriculum Levels};
\srcref{docs/concept/didactic/curriculum-time-axis-and-pacing.md}{Curriculum Time Axis and Pacing}.\\
Prüfungsbasierter Abschluss &
\srcref{docs/concept/skill-graph/assessment-gated-completion.md}{Assessment-Gated Completion};
\srcref{docs/concept/runtime-workflows/exam-preparation-planning.md}{Exam Preparation Planning}.\\
Herkunft und Materialien &
\srcref{docs/concept/skill-graph/source-and-resource-links.md}{Source and Resource Links};
\srcref{docs/concept/skill-graph/content-integration.md}{Content Integration};
\srcref{docs/concept/skill-graph/human-readable-source-rationales.md}{Human-Readable Source Rationales}.\\
Bilder und Zielbücher &
\srcref{docs/concept/skill-graph/atomic-goal-visualizations.md}{Atomic Goal Visualizations};
\srcref{docs/concept/skill-graph/learning-goal-book-and-evidence-review-pipeline.md}{Learning-Goal Book and Evidence-Review Pipeline}.\\
Validierung und QA &
\srcref{docs/qa-ci/graph-validation-rules.md}{Graph Validation Rules};
\srcref{docs/qa-ci/curriculum-quality-maturity-and-routes.md}{Curriculum Quality Maturity and Routes};
\srcref{docs/qa-ci/semantic-atomicity-review.md}{Semantic Atomicity Review}.\\
Austausch und Rechte &
\srcref{docs/concept/skill-graph/dual-curriculum-package-releases.md}{Dual Curriculum Package Releases};
\srcref{docs/concept/skill-graph/mem-fwu-roundtrip-plan.md}{MEM/FWU Roundtrip Plan};
\srcref{LICENSING.md}{LICENSING.md}.\\
\bottomrule
\end{longtable}

Die sieben didaktischen Graphabbildungen stammen aus
\href{https://github.com/enpasos/skillpilot/tree/main/docs/concept/skill-graph/assets/graph-definition}{SkillPilot: Graph Definition, Abbildungsbestand}
und werden unverändert mit deutschen Erläuterungen verwendet.
Die Diagramme auf dem Titelblatt sowie zu Levelplanung und Freistellung
entstehen direkt in TikZ. Maßgeblich sind die Formeln und Bildlegenden;
die Bilder behaupten keine zusätzlichen Kanten oder Anforderungen.

Technischer Text, Formelspezifikation und Build-Werkzeuge stehen unter
Apache-2.0; die didaktischen Illustrationen unter CC BY 4.0, soweit die
entsprechenden Rechte vorliegen. Lizenztexte und Abgrenzung:
\srcref{LICENSE}{LICENSE},
\srcref{LICENSES/CC-BY-4.0.txt}{LICENSES/CC-BY-4.0.txt} und
\srcref{LICENSING.md}{LICENSING.md}.
Quellenhinweis für die Illustrationen: SkillPilot, enpasos -- Enterprise
Patterns \& Solutions GmbH, Graph Definition, Fassung vom 29. Oktober 2026.
Zusammenführung und deutsche Erläuterungen: diese Dokumentfassung.

\section{Dokument bauen und prüfen}
Die Hauptdatei \code{skillpilot-graph-definition.tex} bindet die Kapitel
aus \code{latex/} und die PNG-Abbildungen aus
\code{assets/graph-definition/} ein. Die neuen Diagramme sind TikZ-Vektorgrafiken.

Vom Repositorywurzelverzeichnis aus:
\begin{quote}\small\ttfamily
python3 docs/concept/skill-graph/latex/build-document.py\par
python3 docs/concept/skill-graph/latex/check-assessment-model.py
\end{quote}
Benötigt werden Python 3 und LuaLaTeX mit deutscher Sprachunterstützung
und den Präambelpaketen. Hilfsdateien liegen unter
\code{tmp/skill-graph-document/}, das PDF neben der Hauptdatei.
Der Build löst Verweise in mehreren Läufen auf und bindet Eingabedateien
und PDF durch SHA-256 im Build-Manifest samt Datum und Werkzeugversion.
Über unterschiedliche TeX-Versionen hinweg ist Byteidentität nicht garantiert.

Zum Abschluss gehören ein fehlerfreier LaTeX-Lauf, aufgelöste Verweise,
die Prüfung auf überbreite beziehungsweise überhohe Boxen und eine
visuelle Sichtung des erzeugten PDF. Die Formelprüfung ergänzt diese
Dokumentprüfungen, ersetzt aber weder Runtime-Tests noch unabhängige
fachliche und empirische Validierung.
